\documentclass[10pt,a4paper]{article}

% --- Encodage et langue ---
\usepackage{fontspec}
\usepackage{polyglossia}
\setmainlanguage{french}

% --- Mise en page ---
\usepackage[a4paper, margin=2.5cm]{geometry}
\usepackage{setspace}
\onehalfspacing

% --- Typographie ---
\usepackage{microtype}
\usepackage{parskip}
\setlength{\parindent}{0pt}

% --- Couleurs et boites ---
\usepackage{xcolor}
\definecolor{bleuSGN}{RGB}{0,0,0}
\definecolor{vertSGN}{RGB}{15,100,80}
\definecolor{rougeSGN}{RGB}{180,40,40}
\definecolor{grisbg}{RGB}{248,248,250}
\definecolor{grisrule}{RGB}{200,200,210}
\definecolor{ambre}{RGB}{180,110,10}
\usepackage{tcolorbox}
\tcbuselibrary{skins,breakable,listings}

% --- Code source ---
\usepackage{listings}

\lstdefinelanguage{SGN}{
  morekeywords={ANNEE,NIVEAU,ETUDIANT,MATRICULE,NOM,PRENOM,SEMESTRE,MODULE,COEF,NOTE,L1,L2,L3,S1,S2,S3,S4,S5,S6},
  morecomment=[l]{--},
  morestring=[b]",
  sensitive=true
}

\lstset{
  basicstyle=\ttfamily\small,
  breaklines=true,
  frame=single,
  backgroundcolor=\color{grisbg},
  rulecolor=\color{grisrule},
  numbers=left,
  numberstyle=\tiny\color{gray},
  numbersep=8pt,
  xleftmargin=2em,
  framexleftmargin=2em,
  showstringspaces=false,
  tabsize=4,
  captionpos=b
}

\lstdefinestyle{styleC}{language=C,
  keywordstyle=\color{bleuSGN}\bfseries,
  commentstyle=\color{gray}\itshape,
  stringstyle=\color{rougeSGN},
  morekeywords={Programme,Niveau,Etudiant,Semestre,Module,ErreurSGN,
                CodeNiveau,CodeSemestre,NIV_L1,NIV_L2,NIV_L3,
                SEM_S1,SEM_S2,SEM_S3,SEM_S4,SEM_S5,SEM_S6}
}

\lstdefinestyle{styleFlex}{language=C,
  keywordstyle=\color{vertSGN}\bfseries,
  commentstyle=\color{gray}\itshape,
  stringstyle=\color{rougeSGN}
}

\lstdefinestyle{styleBison}{language=C,
  keywordstyle=\color{ambre}\bfseries,
  commentstyle=\color{gray}\itshape,
  stringstyle=\color{rougeSGN}
}

\lstdefinestyle{styleJSON}{
  language={},
  basicstyle=\ttfamily\small,
  stringstyle=\color{vertSGN},
  keywordstyle=\color{bleuSGN},
  morestring=[b]",
  showstringspaces=false
}

% --- Tableaux ---
\usepackage{amssymb}
\usepackage{tabularx}
\usepackage{booktabs}
\usepackage{array}
\usepackage{multirow}
\renewcommand{\arraystretch}{1.35}

% --- Graphiques / tikz ---
\usepackage{tikz}
\usetikzlibrary{shapes.geometric,arrows.meta,positioning,fit,backgrounds,calc}
\usepackage{float}  % Ajoute cette ligne dans ton préambule

\tikzset{
  boite/.style={rectangle, rounded corners=4pt, draw=bleuSGN, fill=blue!6,
                text width=2.5cm, align=center, minimum height=0.7cm, font=\small},
  boiteC/.style={boite, fill=green!8, draw=vertSGN},
  boiteP/.style={boite, fill=orange!8, draw=ambre},
  boiteE/.style={boite, fill=red!6, draw=rougeSGN},
  fleche/.style={-{Stealth[length=2.5mm]}, thick, draw=gray!70},
  etiq/.style={font=\footnotesize, text=gray!80}
}

% --- En-tetes et pieds de page ---
\usepackage{fancyhdr}
\setlength{\headheight}{15pt}
\pagestyle{fancy}
\fancyhf{}
\fancyhead[L]{\small\textit{Rapport SGN — Compilation}}
\fancyhead[R]{\small\textit{2025--2026}}
\fancyfoot[C]{\thepage}
\renewcommand{\headrulewidth}{0.4pt}

% --- Liens ---
\usepackage[colorlinks=true,linkcolor=bleuSGN,urlcolor=bleuSGN,citecolor=bleuSGN]{hyperref}

% --- Titres ---
\usepackage{titlesec}
\titleformat{\section}[block]
  {\large\bfseries\color{bleuSGN}}
  {\thesection.}{0.5em}{}
  [\vspace{-0.3em}\rule{\linewidth}{0.8pt}]
\titleformat{\subsection}[block]
  {\normalsize\bfseries\color{gray!70!black}}
  {\thesubsection.}{0.5em}{}
\titleformat{\subsubsection}[block]
  {\small\bfseries\color{gray!60!black}}
  {\thesubsubsection.}{0.5em}{}

% --- Boites semantiques (tcbset) ---
\tcbset{
  boxbase/.style={breakable,enhanced,fonttitle=\small\bfseries,
                  left=6pt,right=6pt,top=4pt,bottom=4pt},
  infobox/.style={boxbase,colback=blue!5!white,colframe=bleuSGN!70!white},
  codebox/.style={boxbase,colback=grisbg,colframe=grisrule},
  warnbox/.style={boxbase,colback=red!4!white,colframe=rougeSGN!70!white},
  okbox/.style={boxbase,colback=green!5!white,colframe=vertSGN!70!white}
}

% ============================================================
\begin{document}

% ---- PAGE DE GARDE ----
\begin{titlepage}
  \centering

  % ── Logos ──────────────────────────────────────────────
  \begin{minipage}{0.35\linewidth}
    \centering
    \includegraphics[width=4cm]{/Users/Apple/Downloads/IMG_UGB.jpeg}
  \end{minipage}
  \hfill
  \begin{minipage}{0.35\linewidth}
    \centering
    \includegraphics[width=4cm]{/Users/Apple/Downloads/images.png}
  \end{minipage}

  \vspace{1.2cm}

  % ── Institution ────────────────────────────────────────
  {\large\bfseries Université Gaston Berger de Saint-Louis\\[0.2em]
                   Institut Polytechnique de Saint-Louis\\[0.2em]
                   GeIT 2\par}

  \vspace{0.6cm}

  % ── Bandeau titre ──────────────────────────────────────
  \noindent\textcolor{bleuSGN}{\rule{\linewidth}{1.5pt}}\\[2pt]
  \vspace{0.3cm}

  {\large\bfseries Système de Gestion des Notes (SGN) :\\
  Analyse Lexicale et Syntaxique avec Flex, Bison\par}

  \vspace{0.2cm}
  \noindent\textcolor{bleuSGN}{\rule{\linewidth}{1.5pt}}

  \vspace{2.5cm}

  % ── Auteurs ────────────────────────────────────────────
  \begin{flushleft}
    {\normalsize Réalisés par :\\[0.4em]
    \textbf{Papa Mbaye BA}\\[0.2em]
    \textbf{Gilbert Ngor Thiomby DIOUF}\par}
  \end{flushleft}

  \vspace{1.0cm}

  % ── Encadrant ──────────────────────────────────────────
  \begin{flushright}
    {\normalsize Enseignant :\\[0.4em]
    \textbf{Dr El Hadji Malick NDOYE}\par}
  \end{flushright}

  \vfill

  % ── Date ───────────────────────────────────────────────
  \begin{flushleft}
    {\normalsize Date de Soumission : 15 Juin 2026\par}
  \end{flushleft}

\end{titlepage}

\tableofcontents
\newpage

% ============================================================
\section{Introduction}

\subsection{Contexte et motivation}

Ce projet s'inscrit dans le module de \textbf{Compilation}.
Il consiste à concevoir et implémenter un \textbf{Système de Gestion des Notes (SGN)}
pour les trois premières années de Licence (L1, L2, L3), en mettant en œuvre l'ensemble
de la chaîne de traitement d'un langage dédié : de la grammaire formelle jusqu'à l'interface
graphique.

\subsection{Objectifs pédagogiques}

Le projet couvre les compétences suivantes :

\begin{itemize}
  \item Conception d'une \textbf{grammaire formelle BNF} pour un langage de description de données
  \item Implémentation d'un \textbf{analyseur lexical} avec Flex 2.6
  \item Implémentation d'un \textbf{analyseur syntaxique et sémantique} avec Bison 3.x (LALR-1)
  \item Programmation en \textbf{C (C99)} pour les structures de données et la sérialisation JSON
  \item Développement d'une \textbf{interface graphique} Python/PyQt6 communiquant avec le binaire C via \texttt{subprocess}
\end{itemize}

\subsection{Flux global d'exécution}

\begin{figure}[h]
\centering
\begin{tikzpicture}[node distance=0.8cm]
 % Rangée 1 : pipeline C
  \node[boite]  (sgn)   {Fichier \texttt{.sgn}};
  \node[boiteC, right=1.0cm of sgn]  (lex)   {Flex\\(lexer.l)};
  \node[boiteC, right=1.0cm of lex]  (bison) {Bison\\(parser.y)};
  \node[boiteC, right=1.0cm of bison](main)  {main.c\\(symboles.h)};
  \node[boiteP, right=1.0cm of main] (json)  {JSON\\stdout};

  % Rangée 2 : côté Python (alignée sous json)
  \node[boiteP, below=1.0cm of json] (gui)  {Interface\\Python};
  \node[boite,  left=1.0cm of gui]   (user) {Résultats\\utilisateur};

  % Flèches rangée 1
  \draw[fleche] (sgn)   -- (lex);
  \draw[fleche] (lex)   -- node[above,font=\tiny\itshape]{tokens}   (bison);
  \draw[fleche] (bison) -- node[above,font=\tiny\itshape]{AST/sém.} (main);
  \draw[fleche] (main)  -- (json);

  % Coude JSON → GUI
  \draw[fleche] (json) -- node[right,font=\tiny\itshape]{subprocess} (gui);
  \draw[fleche] (gui)  -- (user);
\end{tikzpicture}
\caption{Flux global d'exécution du système SGN}
\end{figure}

Le fichier source \texttt{.sgn} est d'abord tokenisé par \textbf{Flex}, qui transmet les
tokens à \textbf{Bison}. Ce dernier construit la structure de données en mémoire, effectue
les vérifications sémantiques et déclenche le calcul des moyennes. Le binaire \texttt{sgn\_parser}
émet alors un flux \textbf{JSON} sur la sortie standard, lu par l'interface Python via
\texttt{subprocess.run()}.

% ============================================================
\section{Spécification du langage \texttt{.sgn}}

\subsection{Description du domaine}

Le système gère les notes d'étudiants sur trois niveaux d'études, chacun couvrant deux semestres :

\begin{center}
\begin{tabularx}{0.85\linewidth}{>{\bfseries}l l l l}
\toprule
Niveau & Désignation & Semestres & Modules types \\
\midrule
L1 & Première année de Licence  & S1, S2 & Algorithmique, Maths, Architecture \\
L2 & Deuxième année de Licence  & S3, S4 & POO, BDD, Systèmes, Compilation    \\
L3 & Troisième année de Licence & S5, S6 & IA, Génie Logiciel, Réseaux        \\
\bottomrule
\end{tabularx}
\end{center}

Chaque étudiant possède un \textbf{matricule unique}, un \textbf{nom}, un \textbf{prénom}
(optionnel), et une liste de modules par semestre, chacun associé à un \textbf{coefficient}
et une \textbf{note} sur 20.

\subsection{Structure d'un fichier \texttt{.sgn}}

\begin{lstlisting}[language=SGN, caption={Extrait du fichier \texttt{test\_valide.sgn}}]
-- Fichier notes : Annee academique 2025-2026
ANNEE "2025-2026"

NIVEAU L1 {
    ETUDIANT {
        MATRICULE : "L1-2025-001"
        NOM : "DIALLO Mamadou"
        PRENOM : "Mamadou"
        SEMESTRE S1 {
            MODULE "Algorithmique" COEF 3 NOTE 14.5
            MODULE "Mathematiques" COEF 4 NOTE 12.0
            MODULE "Architecture"  COEF 2 NOTE 16.0
        }
        SEMESTRE S2 {
            MODULE "Programmation C" COEF 3 NOTE 13.5
            MODULE "Logique"         COEF 3 NOTE 11.0
        }
    }
}
\end{lstlisting}

\subsection{Grammaire formelle (BNF)}

La grammaire complète du langage \texttt{.sgn} est définie en BNF comme suit :

\begin{tcolorbox}[codebox, title={Grammaire BNF du langage .sgn}]
\begin{lstlisting}[numbers=none,frame=none,backgroundcolor=\color{white}]
<programme>       ::= ANNEE STRING <liste_niveaux>
<liste_niveaux>   ::= <niveau> | <liste_niveaux> <niveau>
<niveau>          ::= NIVEAU <id_niveau> '{' <liste_etudiants> '}'
<id_niveau>       ::= L1 | L2 | L3
<liste_etudiants> ::= <etudiant> | <liste_etudiants> <etudiant>
<etudiant>        ::= ETUDIANT '{' <champs_etudiant>
                        <liste_semestres> '}'
<champs_etudiant> ::= MATRICULE ':' STRING NOM ':' STRING
                        [ PRENOM ':' STRING ]
<liste_semestres> ::= <semestre> | <liste_semestres> <semestre>
<semestre>        ::= SEMESTRE <id_sem> '{' <liste_modules> '}'
<id_sem>          ::= S1 | S2 | S3 | S4 | S5 | S6
<liste_modules>   ::= <module> | <liste_modules> <module>
<module>          ::= MODULE STRING COEF ENTIER NOTE REEL
<STRING>          ::= '"' [^"]* '"'
<REEL>            ::= [0-9]+ ('.' [0-9]+)?
<ENTIER>          ::= [0-9]+
\end{lstlisting}
\end{tcolorbox}

\subsection{Tableau des tokens lexicaux}

\begin{center}
\begin{tabularx}{0.95\linewidth}{>{\ttfamily}l l l}
\toprule
\normalfont\textbf{Token} & \textbf{Description} & \textbf{Exemple} \\
\midrule
ANNEE      & Mot-clé de début de fichier    & \texttt{ANNEE}             \\
NIVEAU     & Déclaration d'un niveau        & \texttt{NIVEAU}            \\
ETUDIANT   & Bloc étudiant                  & \texttt{ETUDIANT}          \\
MATRICULE  & Identifiant matricule          & \texttt{MATRICULE}         \\
NOM / PRENOM & Champs nominaux             & \texttt{NOM}, \texttt{PRENOM} \\
SEMESTRE   & Déclaration d'un semestre      & \texttt{SEMESTRE}          \\
MODULE     & Déclaration d'un module        & \texttt{MODULE}            \\
COEF       & Coefficient du module          & \texttt{COEF}              \\
NOTE       & Note obtenue                   & \texttt{NOTE}              \\
L1/L2/L3   & Identifiants de niveau         & \texttt{L1}, \texttt{L3}   \\
S1..S6     & Identifiants de semestre       & \texttt{S1}, \texttt{S4}   \\
REEL       & Nombre réel                    & \texttt{14.5}, \texttt{9.0}\\
ENTIER     & Nombre entier                  & \texttt{3}, \texttt{1}     \\
STRING\_TOK & Chaîne entre guillemets       & \texttt{"Algorithmique"}   \\
\bottomrule
\end{tabularx}
\end{center}

% ============================================================
\section{Analyseur lexical --- \texttt{lexer.l} (Flex)}

\subsection{Structure du fichier}


Le fichier \texttt{lexer.l} est organisé en trois sections selon la convention Flex :

\begin{figure}[h]
\centering
\begin{tikzpicture}[node distance=0.3cm]
  \node[boiteC,text width=12cm,minimum height=0.8cm]
    (s1) {\textbf{Section 1 — Prologue C} : \texttt{\%\{ \#include "parser.tab.h", "symboles.h" \%\}}};
  \node[boiteC,text width=12cm,minimum height=0.8cm,below=of s1]
    (s2) {\textbf{Section 2 — Définitions} : \texttt{\%option yylineno}, macros \texttt{BLANC}, \texttt{NOMBRE}, \texttt{STR}, \texttt{COMMENT}};
  \node[boiteC,text width=12cm,minimum height=0.8cm,below=of s2]
    (s3) {\textbf{Section 3 — Règles} : 19 tokens reconnus + gestion erreurs via \texttt{ajouter\_erreur()}};
  \draw[fleche] (s1) -- (s2);
  \draw[fleche] (s2) -- (s3);
\end{tikzpicture}
\caption{Structure du fichier \texttt{lexer.l}}
\end{figure}

Le fichier \texttt{lexer.l} est le fichier source de l'analyseur lexical écrit pour Flex.
Tout fichier Flex est obligatoirement divisé en trois sections séparées par \texttt{\%\%}.
La première, le \textbf{prologue C}, contient du code C pur encadré par \texttt{\%\{} et
\texttt{\%\}}, copié tel quel en tête du fichier généré : on y place les \texttt{\#include}
de \texttt{parser.tab.h} (qui fournit les numéros de tokens générés par Bison) et de
\texttt{symboles.h} (qui donne accès aux structures de données et à
\texttt{ajouter\_erreur()}). La deuxième section contient les \textbf{définitions} :
l'option \texttt{\%option yylineno} demande à Flex de compter automatiquement les lignes
dans \texttt{yylineno}, et les macros (\texttt{BLANC}, \texttt{NOMBRE}, \texttt{STR},
\texttt{COMMENT}) sont des abréviations d'expressions régulières réutilisables dans la
section suivante.

La troisième section, le \textbf{cœur du lexer}, contient les règles : chacune associe
un motif (expression régulière) à une action (code C). Les 19 règles couvrent les
mots-clés (\texttt{ANNEE}, \texttt{NIVEAU}\ldots), les identifiants de niveau et de
semestre (\texttt{L1}, \texttt{S1}\ldots), les littéraux (\texttt{STRING}, \texttt{REEL},
\texttt{ENTIER}) et les séparateurs (\texttt{\{}, \texttt{\}}, \texttt{:}). Les deux
dernières règles sont dédiées à la gestion des erreurs : elles capturent tout caractère
ou chaîne non reconnu et appellent \texttt{ajouter\_erreur("LEX", yylineno, msg)} pour
enregistrer l'erreur sans interrompre le parsing. Flex lit le fichier \texttt{.sgn}
caractère par caractère, applique ces règles et transmet les tokens reconnus à Bison.


\subsection{Options et définitions de macros}

\begin{lstlisting}[style=styleFlex, caption={Options et macros dans \texttt{lexer.l}}]
%option noyywrap yylineno warn noinput nounput

BLANC    [ \t\r]+
NL       \n
NOMBRE   [0-9]+(\.[0-9]+)?
STR      \"[^\"\n]*\"
COMMENT  "--"[^\n]*
\end{lstlisting}

L'option \texttt{yylineno} délègue à Flex le suivi automatique du numéro de ligne,
indispensable pour la localisation des erreurs. L'option \texttt{noyywrap} supprime
la redéfinition de \texttt{yywrap()}.

\subsection{Règles de reconnaissance}

\begin{lstlisting}[style=styleFlex, caption={Extrait des règles lexicales --- mots-clés et littéraux}]
{COMMENT}   { /* commentaire ignore */ }
{BLANC}     { /* blancs ignores      */ }
{NL}        { /* saut de ligne       */ }

"ANNEE"     { return ANNEE;        }
"NIVEAU"    { return NIVEAU_KW;    }
"ETUDIANT"  { return ETUDIANT_KW;  }
"MODULE"    { return MODULE_KW;    }
"COEF"      { return COEF_KW;      }
"NOTE"      { return NOTE_KW;      }
"L1"        { return NIV_L1_TOK;  }
"S1"        { return SEM_S1_TOK;  }

{NOMBRE}    {
    if (strchr(yytext, '.')) {
        yylval.reel   = atof(yytext); return REEL;
    } else {
        yylval.entier = atoi(yytext); return ENTIER;
    }
}
{STR}       { yylval.chaine = strip_quotes(yytext); return STRING_TOK; }
\end{lstlisting}

La règle \texttt{NOMBRE} dispatche en deux tokens distincts : \textbf{REEL} si le texte
contient un point décimal, \textbf{ENTIER} sinon. La fonction \texttt{strip\_quotes()}
retire les guillemets et retourne une chaîne allouée dynamiquement avec \texttt{malloc()},
libérée dans les actions Bison.

\subsection{Gestion des erreurs lexicales}

\begin{lstlisting}[style=styleFlex, caption={Collecte des erreurs lexicales dans \texttt{g\_erreurs[]}}]
\"[^\"\n]*\n {
    char msg[MAX_STR];
    snprintf(msg, MAX_STR,
        "chaine non terminee a la ligne %d", yylineno);
    ajouter_erreur("LEX", yylineno, msg);
}
.  {
    char msg[MAX_STR];
    snprintf(msg, MAX_STR,
        "caractere illegal '%s' (ASCII %d)",
        yytext, (unsigned char)yytext[0]);
    ajouter_erreur("LEX", yylineno, msg);
}
\end{lstlisting}

Toute erreur lexicale appelle \texttt{ajouter\_erreur("LEX", ...)} qui l'enregistre
dans le tableau global \texttt{g\_erreurs[]} et la signale sur \texttt{stderr}.
Le parsing peut se poursuivre malgré l'erreur.

% ============================================================
\section{Analyseur syntaxique et sémantique --- \texttt{parser.y} (Bison)}

\subsection{Déclarations et union de valeurs}

\begin{lstlisting}[style=styleBison, caption={Déclarations Bison --- union et tokens}]
%define parse.error verbose

%union { double reel; int entier; char *chaine; }

%token ANNEE NIVEAU_KW ETUDIANT_KW MATRICULE_KW NOM_KW PRENOM_KW
%token SEMESTRE_KW MODULE_KW COEF_KW NOTE_KW
%token NIV_L1_TOK NIV_L2_TOK NIV_L3_TOK
%token SEM_S1_TOK SEM_S2_TOK SEM_S3_TOK SEM_S4_TOK SEM_S5_TOK SEM_S6_TOK
%token <reel>   REEL
%token <entier> ENTIER
%token <chaine> STRING_TOK
%type  <entier> id_niveau id_sem
\end{lstlisting}

L'union \texttt{\%union} définit les trois types de valeurs sémantiques possibles.
La directive \texttt{\%define parse.error verbose} produit des messages d'erreur
syntaxique détaillés, intégrés dans le JSON de sortie.

\subsection{Variables de contexte}

Quatre variables statiques maintiennent le contexte de construction pendant le parsing :

\begin{lstlisting}[style=styleBison, caption={Variables de contexte dans \texttt{parser.y}}]
static Niveau    *ctx_niv      = NULL;  /* niveau courant    */
static Etudiant  *ctx_etu      = NULL;  /* etudiant courant  */
static Semestre  *ctx_sem      = NULL;  /* semestre courant  */
static CodeNiveau ctx_code_niv = 0;     /* code du niveau    */
\end{lstlisting}

\subsection{Règles de production principales}

\begin{lstlisting}[style=styleBison, caption={Règle \texttt{programme} et règle \texttt{module\_item}}]
programme
    : ANNEE STRING_TOK liste_niveaux
        {
            strncpy(g_programme.annee, $2, MAX_STR - 1);
            free($2);
            programme_calcule_tout(&g_programme);
        }
    ;

module_item
    : MODULE_KW STRING_TOK COEF_KW ENTIER NOTE_KW REEL
        {
            char msg[MAX_STR];
            /* Verification 1 : note dans [0,20] */
            if ($6 < 0.0 || $6 > 20.0) {
                snprintf(msg, MAX_STR,
                    "note %.2f hors [0,20] pour \"%s\"", $6, $2);
                ajouter_erreur("SEM", yylineno, msg);
            }
            /* Verification 2 : coefficient >= 1 */
            if ($4 < 1) {
                snprintf(msg, MAX_STR,
                    "coefficient %d invalide pour \"%s\"", $4, $2);
                ajouter_erreur("SEM", yylineno, msg);
            }
            if (ctx_sem) semestre_add_module(ctx_sem, $2, $4, $6);
            free($2);
        }
    ;
\end{lstlisting}

\subsection{Les quatre vérifications sémantiques}

\begin{center}
\begin{tabularx}{\linewidth}{>{\bfseries}l l X}
\toprule
\# & Condition vérifiée & Action en cas d'erreur \\
\midrule
1 & Note $\in [0.0,\ 20.0]$ & \texttt{ajouter\_erreur("SEM", ...)} + parsing continue \\
2 & Coefficient $\geq 1$     & \texttt{ajouter\_erreur("SEM", ...)} + parsing continue \\
3 & Matricule unique par niveau & \texttt{ajouter\_erreur("SEM", ...)} + \texttt{ctx\_etu = NULL} \\
4 & Semestre cohérent avec niveau & \texttt{ajouter\_erreur("SEM", ...)} + semestre ignoré \\
\bottomrule
\end{tabularx}
\end{center}

\vspace{0.5em}
La vérification~4 repose sur la fonction \texttt{semestre\_valide\_pour\_niveau()} définie
dans \texttt{main.c} : L1 n'accepte que S1/S2, L2 que S3/S4, L3 que S5/S6.

\subsection{Organigramme du traitement d'une erreur sémantique}

\begin{figure}[h]
\centering
\begin{tikzpicture}[node distance=0.8cm and 1.5cm]
  \node[boiteC] (det) {Détection\\dans action Bison};
  \node[boiteE,right=2cm of det] (add) {\texttt{ajouter\_erreur()}\\stocke dans \texttt{g\_erreurs[]}};
  \node[boiteC,right=2cm of add] (cont) {Parsing\\continue};
  \node[boiteP,below=of add] (err) {Émission sur\\stderr};
  \node[boiteP,below=of cont] (json) {Inclus dans\\JSON de sortie};

  \draw[fleche] (det)  -- (add);
  \draw[fleche] (add)  -- (cont);
  \draw[fleche] (add)  -- (err);
  \draw[fleche] (cont) -- (json);
\end{tikzpicture}
\caption{Flux de traitement d'une erreur sémantique}
\end{figure}

La fonction \texttt{yyerror()} collecte de même les erreurs syntaxiques :

\begin{lstlisting}[style=styleBison, caption={Collecte unifiée dans \texttt{yyerror()}}]
void yyerror(const char *msg)
{
    ajouter_erreur("SYN", yylineno, msg);
}
\end{lstlisting}

% ============================================================
\section{Implémentation C --- \texttt{main.c} et \texttt{symboles.h}}

\subsection{Architecture des fichiers source}

\begin{figure}[h]
\centering
\begin{tikzpicture}[node distance=0.6cm and 1.2cm]
  \node[boiteC,text width=3.5cm] (sh) {\textbf{symboles.h}\\structures + prototypes\\(déclarations uniquement)};
  \node[boite,text width=3cm,above left=1.2cm and 0.8cm of sh]  (ly) {lexer.l\\(Flex)};
  \node[boite,text width=3cm,above=1.2cm of sh]       (py) {parser.y\\(Bison)};
  \node[boite,text width=3cm,above right=1.2cm and 0.8cm of sh] (mc) {main.c\\(implémentations)};
  \node[boiteP,text width=3cm,below=1.2cm of sh]      (mk) {Makefile\\(compilation)};

  \draw[fleche] (sh) -- (ly);
  \draw[fleche] (sh) -- (py);
  \draw[fleche] (sh) -- (mc);
  \draw[fleche] (mk) -- (sh);
\end{tikzpicture}
\caption{Dépendances entre les fichiers source}
\end{figure}

Le fichier \texttt{symboles.h} ne contient \textbf{aucune implémentation} : uniquement
les déclarations de structures, les constantes, les \texttt{extern} des variables globales
et les prototypes de fonctions. Toutes les implémentations résident dans \texttt{main.c}.

\subsection{Structures de données}

\begin{lstlisting}[style=styleC, caption={Hiérarchie des structures dans \texttt{symboles.h}}]
typedef struct {
    char   nom[MAX_STR]; int coef; double note;
} Module;

typedef struct {
    CodeSemestre id;
    Module       modules[MAX_MOD];  /* MAX_MOD = 64 */
    int          nb_modules;
    double       moyenne;
} Semestre;

typedef struct {
    char     matricule[MAX_STR], nom[MAX_STR];
    Semestre semestres[MAX_SEM]; /* MAX_SEM = 6  */
    int      nb_semestres;
    double   moyenne_annuelle;
    char     mention[32], decision[16];
    int      rang;
} Etudiant;

typedef struct {
    CodeNiveau id;
    Etudiant   etudiants[MAX_ETU]; /* MAX_ETU = 256 */
    int        nb_etudiants;
} Niveau;

typedef struct {
    char   annee[MAX_STR];
    Niveau niveaux[MAX_NIV];       /* MAX_NIV = 3   */
    int    nb_niveaux;
} Programme;

typedef struct {
    int  ligne; char type[16]; char message[MAX_STR];
} ErreurSGN;
\end{lstlisting}

Le choix des \textbf{tableaux statiques} (plutôt que l'allocation dynamique) simplifie
la gestion mémoire : aucun \texttt{free()} nécessaire sur les structures, les limites
\texttt{MAX\_*} sont suffisamment larges pour le contexte universitaire visé.

\subsection{Variables globales}

\begin{lstlisting}[style=styleC, caption={Variables globales définies dans \texttt{main.c}}]
Programme g_programme;          /* arbre de donnees principal  */
ErreurSGN g_erreurs[MAX_ERR];   /* toutes erreurs LEX/SYN/SEM  */
int       g_nb_erreurs = 0;     /* compteur d'erreurs          */
\end{lstlisting}

Ces variables sont accessibles depuis \texttt{lexer.l} et \texttt{parser.y} grâce
aux déclarations \texttt{extern} dans \texttt{symboles.h}. Ce choix est justifié par
la nature du pipeline Flex/Bison : les actions lexicales et syntaxiques doivent pouvoir
alimenter la structure globale sans passer de pointeurs à travers la pile LALR.

\subsection{Catalogue des fonctions implémentées}

\begin{center}
\begin{tabularx}{\linewidth}{>{\ttfamily\small}l X}
\toprule
\normalfont\textbf{Fonction} & \textbf{Rôle} \\
\midrule
programme\_init(p)           & Initialise la structure \texttt{Programme} à zéro (\texttt{memset}) \\
programme\_add\_niveau(p, code) & Recherche ou crée le niveau \texttt{code} dans \texttt{p} \\
niveau\_find\_etudiant(niv, mat) & Recherche linéaire par matricule dans le niveau \\
niveau\_add\_etudiant(niv, mat, nom) & Alloue un nouvel étudiant dans le niveau \\
etudiant\_add\_semestre(e, id) & Crée un semestre pour l'étudiant courant \\
semestre\_add\_module(s, nom, coef, note) & Ajoute un module au semestre courant \\
programme\_calcule\_tout(p) & Orchestre les calculs post-parsing \\
semestre\_valide\_pour\_niveau(sem, niv) & Vérifie la cohérence semestre/niveau \\
ajouter\_erreur(type, ligne, msg) & Enregistre toute erreur LEX/SYN/SEM \\
programme\_to\_json(p, out) & Sérialise la structure complète en JSON \\
\bottomrule
\end{tabularx}
\end{center}

\subsection{Calculs --- formules du sujet}

\subsubsection{Moyenne pondérée d'un semestre (formule 1)}

\begin{lstlisting}[style=styleC, caption={Implémentation de la formule (1)}]
static double semestre_moyenne(Semestre *s)
{
    double somme_pond = 0.0, somme_coef = 0.0;
    int i;
    for (i = 0; i < s->nb_modules; i++) {
        somme_pond += s->modules[i].coef * s->modules[i].note;
        somme_coef += s->modules[i].coef;
    }
    return somme_coef > 0.0 ? somme_pond / somme_coef : 0.0;
}
\end{lstlisting}

$$\bar{M}_s = \frac{\displaystyle\sum_{i=1}^{n} c_i \times n_i}{\displaystyle\sum_{i=1}^{n} c_i}$$

\subsubsection{Moyenne annuelle (formule 2)}

\begin{lstlisting}[style=styleC, caption={Implémentation de la formule (2)}]
static double etudiant_moyenne_annuelle(Etudiant *e)
{
    double total = 0.0;
    int i;
    if (e->nb_semestres == 0) return 0.0;
    for (i = 0; i < e->nb_semestres; i++)
        total += e->semestres[i].moyenne;
    return total / e->nb_semestres;
}
\end{lstlisting}

$$\bar{M}_a = \frac{1}{S}\sum_{s=1}^{S} \bar{M}_s$$

\subsubsection{Décision et mention}

\begin{lstlisting}[style=styleC, caption={Attribution de la mention}]
static void etudiant_set_mention(Etudiant *e)
{
    double m = e->moyenne_annuelle;
    if      (m >= 16.0) { strcpy(e->decision,"Admis"); strcpy(e->mention,"Tres Bien");  }
    else if (m >= 14.0) { strcpy(e->decision,"Admis"); strcpy(e->mention,"Bien");       }
    else if (m >= 12.0) { strcpy(e->decision,"Admis"); strcpy(e->mention,"Assez Bien"); }
    else if (m >= 10.0) { strcpy(e->decision,"Admis"); strcpy(e->mention,"Passable");   }
    else                { strcpy(e->decision,"Ajourne"); strcpy(e->mention,"");          }
}
\end{lstlisting}

\subsubsection{Tri et calcul des rangs}

\textbf{niveau\_calcule\_rangs()} trie le tableau \texttt{etudiants[]} par ordre
décroissant de \texttt{moyenne\_annuelle} via un tri à bulles, puis affecte le rang
\texttt{i+1} à l'étudiant à la position \texttt{i}.

\subsection{Organigramme de \texttt{main()}}

\begin{figure}[H]
\centering
\begin{tikzpicture}[node distance=0.55cm,]
  \node[boiteC] (args)   {Parse arguments\\(\texttt{fichier}, \texttt{--debug})};
  \node[boiteC,below=of args]  (open)   {\texttt{fopen(fichier)}};
  \node[boiteC,below=of open]  (init)   {\texttt{programme\_init()}};
  \node[boiteC,below=of init]  (yyr)    {\texttt{yyrestart(fp)}};
  \node[boiteC,below=of yyr]   (yp)     {\texttt{yyparse()}};
 \node[boiteC, below=of yp, text width=3.2cm, font=\small] (calc) {\texttt{programme\_}\\ \texttt{calculate\_tout()}};
 \node[boiteC, below=of calc, text width=3.2cm, font=\small] (json) {\texttt{programme\_to\_json}\\ \texttt{(stdout)}};
 \node[boiteP,below=of json]  (exit)   {\texttt{EXIT\_SUCCESS} / \texttt{EXIT\_FAILURE}};

  \draw[fleche] (args) -- (open);
  \draw[fleche] (open) -- (init);
  \draw[fleche] (init) -- (yyr);
  \draw[fleche] (yyr)  -- (yp);
  \draw[fleche] (yp)   -- (calc);
  \draw[fleche] (calc) -- (json);
  \draw[fleche] (json) -- (exit);
\end{tikzpicture}
\caption{Organigramme de la fonction \texttt{main()}}
\end{figure}

\subsection{Sérialisation JSON}

La fonction \texttt{programme\_to\_json()} parcourt la hiérarchie complète
\texttt{Programme → Niveau → Etudiant → Semestre → Module} et produit le JSON
conforme au format §3.3.2 du sujet. La fonction auxiliaire \texttt{json\_esc()} échappe
les caractères spéciaux (\texttt{"} et \texttt{\textbackslash}) pour garantir un JSON valide
même si les noms contiennent des guillemets.

\begin{tcolorbox}[okbox, title={Exemple de sortie JSON pour \texttt{test\_valide.sgn}}]
\begin{lstlisting}[style=styleJSON, numbers=none, frame=none, backgroundcolor=\color{green!5}]
{
  "annee": "2025-2026",
  "niveaux": [
    {
      "niveau": "L1",
      "etudiants": [
        {
          "matricule": "L1-2025-001",
          "nom": "DIALLO Mamadou",
          "moyenne_annuelle": 13.64,
          "mention": "Assez Bien",
          "decision": "Admis",
          "rang": 1
        }
      ]
    }
  ],
  "erreurs": []
}
\end{lstlisting}
\end{tcolorbox}

% ============================================================
\section{Compilation --- \texttt{Makefile}}

\subsection{Compatibilité multi-plateformes}

Le Makefile détecte automatiquement le système d'exploitation via \texttt{\$(shell uname)}
et adapte les flags en conséquence :

\begin{lstlisting}[language=bash, caption={Détection OS dans le Makefile}]
OS := $(shell uname)

ifeq ($(OS), Darwin)          # macOS
    BREW_BISON := $(shell brew --prefix bison 2>/dev/null)/bin/bison
    ifeq ($(shell test -x "$(BREW_BISON)" && echo yes), yes)
        BISON = $(BREW_BISON)
        BISON_WARN = -Wall      # bison Homebrew recent
    else
        BISON = bison
        BISON_WARN =            # bison Xcode (vieux, sans -Wall)
    endif
    LDFLAGS = -ll -lm
else                           # Linux (Ubuntu/Debian)
    BISON = bison
    BISON_WARN = -Wall
    LDFLAGS = -lfl -lm
endif
\end{lstlisting}

\subsection{Chaîne de compilation}

\begin{figure}[H]
\centering
\begin{tikzpicture}[node distance=0.6cm and 1.2cm]
  \node[boite]  (ly)  {\texttt{lexer.l}};
  \node[boite,right=2cm of ly]  (py)  {\texttt{parser.y}};
  \node[boite,right=2cm of py]  (mc)  {\texttt{main.c}};

  \node[boiteC,below=of ly]   (lex)   {\texttt{flex}};
  \node[boiteC,below=of py]   (bis)   {\texttt{bison -d}};
  \node[boiteP,below=of bis]  (tab)   {\texttt{parser.tab.h}};

  \node[boiteP,below=of lex]  (lyy)   {\texttt{lex.yy.o}};
  \node[boiteP,below=of tab]  (tabo)  {\texttt{parser.tab.o}};
  \node[boiteP,below=of mc]   (mo)    {\texttt{main.o}};

  \node[boiteE,below=2.2cm of tabo] (bin) {\texttt{sgn\_parser}};

  \draw[fleche] (ly)  -- (lex);
  \draw[fleche] (py)  -- (bis);
  \draw[fleche] (bis) -- (tab);
  \draw[fleche] (lex) -- (lyy);
  \draw[fleche] (tab) -- (lyy);
  \draw[fleche] (tab) -- (tabo);
  \draw[fleche] (mc)  -- (mo);
  \draw[fleche] (lyy)  -- (bin);
  \draw[fleche] (tabo) -- (bin);
  \draw[fleche] (mo)   -- (bin);
\end{tikzpicture}
\caption{Chaîne de compilation du projet SGN}
\end{figure}

% ============================================================
\section{Interface graphique Python}

\subsection{Architecture des fichiers}

L'interface est répartie en trois fichiers Python :

\begin{center}
\begin{tabularx}{0.92\linewidth}{>{\ttfamily}l X}
\toprule
\normalfont\textbf{Fichier} & \textbf{Rôle} \\
\midrule
app.py          & Fenêtre principale PyQt6 : Vue + Contrôleur \\
parser\_bridge.py & Pont subprocess vers \texttt{sgn\_parser}, parse le JSON retourné \\
interface.py    & Constantes de style, widget \texttt{FilePanel} (drag \& drop) \\
\bottomrule
\end{tabularx}
\end{center}

\subsection{Diagramme de séquence : analyse d'un fichier}

\begin{figure}[H]
\centering
\begin{tikzpicture}[node distance=1.0cm and 0.5cm]
  % Acteurs
  \node[boite,text width=2.2cm]  (usr)  {Utilisateur};
  \node[boiteP,text width=2.4cm,right=1.8cm of usr] (app)  {app.py};
  \node[boiteP,text width=2.8cm,right=1.8cm of app] (bri)  {parser\_bridge.py};
  \node[boiteC,text width=2.4cm,right=1.8cm of bri] (bin)  {sgn\_parser};
  
  \draw[dashed, thin, gray!90] (usr.south) -- ([yshift=-6cm]usr.south);
  \draw[dashed, thin, gray!90] (app.south) -- ([yshift=-6cm]app.south);
  \draw[dashed, thin, gray!90] (bri.south) -- ([yshift=-6cm]bri.south);
  \draw[dashed, thin, gray!90] (bin.south) -- ([yshift=-6cm]bin.south);

  % Messages séquentiels (flèches simples vers la droite et gauche)
  \node[below=0.6cm of usr,etiq,anchor=west] (m1) {};
  \draw[fleche] ([yshift=-1.0cm]usr.south) -- node[above,etiq,font=\footnotesize,text=darkgray]{(1) Clic Analyser} ([yshift=-1.0cm]app.south);
  \draw[fleche] ([yshift=-1.8cm]app.south) -- node[above,etiq,font=\footnotesize, text=darkgray]{(2) analyser(chemin)} ([yshift=-1.8cm]bri.south);
  \draw[fleche] ([yshift=-2.6cm]bri.south) -- node[above,etiq,font=\footnotesize,text=darkgray]{(3) subprocess.run()} ([yshift=-2.6cm]bin.south);
  \draw[fleche] ([yshift=-3.6cm]bin.south) -- node[above,etiq,font=\footnotesize, text=darkgray]{(4) JSON stdout} ([yshift=-3.6cm]bri.south);
  \draw[fleche] ([yshift=-4.4cm]bri.south) -- node[above,etiq,font=\footnotesize, text=darkgray]{(5) dict Python} ([yshift=-4.4cm]app.south);
  \draw[fleche] ([yshift=-5.2cm]app.south) -- node[above,etiq,font=\footnotesize, text=darkgray]{(6) Affichage résultats} ([yshift=-5.2cm]usr.south);
\end{tikzpicture}
\caption{Diagramme de séquence : analyse d'un fichier \texttt{.sgn}}
\end{figure}

\subsection{Communication subprocess}

\begin{lstlisting}[language=Python, caption={Extrait de \texttt{parser\_bridge.py} --- appel subprocess}]
res = subprocess.run(
    [bin_path, chemin_sgn],
    capture_output=True, text=True, timeout=15
)
if res.stdout.strip():
    return json.loads(res.stdout)
\end{lstlisting}

La communication entre Python et le binaire C est \textbf{unidirectionnelle} :
le binaire lit le fichier \texttt{.sgn} depuis le disque, écrit le JSON sur
\texttt{stdout} et les messages de débogage sur \texttt{stderr}. Python capture
\texttt{stdout} via \texttt{capture\_output=True} et décode le JSON avec
\texttt{json.loads()}.

\subsection{Fonctionnalités de l'interface}

L'interface offre les fonctionnalités suivantes :

\begin{itemize}
  \item \textbf{Chargement} par bouton Parcourir ou drag \& drop sur le panneau fichier
  \item \textbf{Éditeur intégré} : zone de texte éditable avant analyse, sauvegardée avant chaque appel
  \item \textbf{Onglet Résultats} : tableau PyQt6 avec colonnes Matricule, Nom, Niveau, Semestre, Moyennes, Rang, Mention, Décision
  \item \textbf{Onglet JSON brut} : affichage du JSON complet retourné par \texttt{sgn\_parser}
  \item \textbf{Bandeau d'erreurs} : apparaît uniquement si des erreurs sont détectées, avec type (LEX/SYN/SEM), numéro de ligne et message
  \item \textbf{Export CSV} : tableau des résultats + section erreurs en bas
  \item \textbf{Export TXT} : tableau ASCII avec colonnes centrées + section erreurs
\end{itemize}

% ============================================================
\section{Tests et validation}

\subsection{Stratégie de tests}

Quatre fichiers de test couvrent les cas nominaux et d'erreur :

\begin{center}
\begin{tabularx}{\linewidth}{>{\ttfamily}l l X}
\toprule
\normalfont\textbf{Fichier} & \textbf{Type} & \textbf{Contenu} \\
\midrule
test\_valide.sgn     & Nominal       & 3 étudiants sur 2 niveaux, toutes données correctes \\
test\_erreur\_lex.sgn & Erreur LEX   & Caractère illégal \texttt{\#}, chaîne non terminée \\
test\_erreur\_syn.sgn & Erreur SYN   & Token manquant, accolade non fermée \\
test\_erreur\_sem.sgn & Erreur SEM   & Note 22.0, coef 0, doublon matricule, S3 dans L1 \\
\bottomrule
\end{tabularx}
\end{center}

\subsection{Vérification de la compilation}

\begin{tcolorbox}[okbox, title={Sortie de \texttt{make} sur Ubuntu 24.04}]
\begin{figure}[H]
	\centering
	\includegraphics[width=\textwidth,keepaspectratio]{/Users/Apple/Downloads/Capture d’écran 2026-06-11 à 08.37.03.png}
\end{figure}
\end{tcolorbox}


\subsection{Résultats obtenus}

\subsubsection{Test 1 --- Fichier valide}

\begin{figure}[H]  
	\centering
	\includegraphics[width=\textwidth,keepaspectratio]{/Users/Apple/Downloads/Capture d’écran 2026-06-11 à 08.47.47.png}
	\caption{Test d'un fichier .sgn valide}
\end{figure}

\subsubsection{Test 4 --- Erreurs sémantiques}

\begin{figure}[H]
	\centering
	\includegraphics[width=\textwidth,keepaspectratio]{/Users/Apple/Downloads/Capture d’écran 2026-06-11 à 08.48.12.png}
	\caption{Test d'un fichier .sgn contenant des erreurs sémantiques}
\end{figure}

\textbf{Point clé} : le JSON est \textbf{toujours émis}, même en présence d'erreurs.
Les erreurs apparaissent dans le tableau \texttt{"erreurs":[...]} du JSON, rendant
l'interface Python opérationnelle dans tous les cas.


% ============================================================
\section{Conclusion}

\subsection{Bilan des fonctionnalités}

\begin{center}
\begin{tabularx}{\linewidth}{X >{\centering\arraybackslash}l >{\centering\arraybackslash}l}
\toprule
\textbf{Fonctionnalité} & \textbf{Réalisé} \\
\midrule
Analyseur lexical Flex (tous tokens, erreurs, commentaires)  & \checkmark \\
Analyseur syntaxique Bison (grammaire BNF conforme)          & \checkmark \\
Construction table/structure de données                       & \checkmark \\
Vérifications sémantiques (4 règles)                          & \checkmark \\
Calculs moyennes, rangs, mentions                            & \checkmark \\
Interface graphique PyQt6 (chargement, édition, tableau)       & \checkmark \\
Affichage erreurs coloré avec numéro de ligne                  & \checkmark \\
Export CSV et TXT                                              & \checkmark \\
Rapport technique                                              & \checkmark \\
\bottomrule
\end{tabularx}
\end{center}

\subsection{Difficultés rencontrées}

Les principales difficultés rencontrées ont été les suivantes. La gestion des
variables de contexte Bison (\texttt{ctx\_niv}, \texttt{ctx\_etu}, \texttt{ctx\_sem})
a nécessité une attention particulière : une affectation prématurée ou un oubli de
réinitialisation à \texttt{NULL} provoquait des erreurs de segmentation silencieuses.
La collecte unifiée des erreurs LEX, SYN et SEM dans un seul tableau JSON a
requis la refactorisation de \texttt{yyerror()} et des règles du lexer pour
appeler \texttt{ajouter\_erreur()} de façon cohérente.
La compatibilité macOS/Linux au niveau du Makefile a demandé la détection de
l'OS via \texttt{\$(shell uname)} et l'adaptation des flags \texttt{-Wall} (non supporté
par le bison embarqué dans Xcode) et des bibliothèques de liaison (\texttt{-ll} vs \texttt{-lfl}).

\subsection{Perspectives d'amélioration}

Plusieurs améliorations pourraient être envisagées. L'ajout d'un \textbf{export PDF}
directement depuis l'interface Python 
permettrait une remise de bulletins plus professionnelle. La prise en charge de
\textbf{L4 et L5} nécessiterait uniquement d'étendre les énumérations \texttt{CodeNiveau}
et \texttt{CodeSemestre} et d'ajouter les tokens correspondants dans le lexer.
Enfin, une saisie directe dans la GUI (sans passer par un fichier disque
intermédiaire) améliorerait l'ergonomie pour les utilisateurs avancés.

\end{document}
